%==============================================================
%  abm.tex
%  hasklib — mathematical companion / derivations
%
%  Arithmetic Brownian Motion (Bachelier's model): exact solution
%  by direct integration -- no Ito's lemma or integrating factor
%  needed, since the coefficients are already constant.
%
%  Sibling of gbm.tex and ou.tex, not a dependent of either.
%==============================================================
\documentclass[11pt]{article}

%---------------- packages ----------------
\usepackage[utf8]{inputenc}
\usepackage[T1]{fontenc}
\usepackage{amsmath, amssymb, amsthm}
\usepackage{mathtools}
\usepackage[a4paper, margin=1in]{geometry}
\usepackage{xcolor}
\usepackage{hyperref}
\usepackage{cleveref}
\usepackage{bm}
\usepackage{harvard}   % Harvard author-date referencing (with agsm.bst)

%---------------- notation macros ----------------
% (kept identical to ito_foundations.tex / gbm.tex / ou.tex)
\newcommand{\Prob}{\mathbb{P}}
\newcommand{\Qmeas}{\mathbb{Q}}
\newcommand{\E}{\mathbb{E}}
\newcommand{\Var}{\operatorname{Var}}
\newcommand{\Cov}{\operatorname{Cov}}
\newcommand{\filt}{\mathcal{F}}
\newcommand{\Om}{\Omega}

\newcommand{\dd}{\mathrm{d}}
\newcommand{\dW}{\dd W_t}
\newcommand{\dt}{\dd t}
\newcommand{\R}{\mathbb{R}}

\newcommand{\drift}{a}
\newcommand{\diff}{b}

%---------------- theorem environments ----------------
\theoremstyle{plain}
\newtheorem{theorem}{Theorem}
\newtheorem{proposition}{Proposition}

\theoremstyle{definition}
\newtheorem{definition}{Definition}

\theoremstyle{remark}
\newtheorem{remark}{Remark}

%==============================================================
\title{Arithmetic Brownian Motion: Exact Solution\\
       \large \texttt{hasklib} mathematical companion}
\author{Hubert}
\date{\today}
%==============================================================

\begin{document}
\maketitle

\begin{abstract}
This note derives the exact terminal law of arithmetic Brownian
motion (ABM), the process underlying \citeasnoun{shreve2004}'s
classical treatment and, historically, Bachelier's original 1900
model of a stock price. Unlike GBM and OU, neither It\^o's lemma nor
an integrating factor is needed: both coefficients are already
constant, so the SDE integrates directly.
\end{abstract}

%--------------------------------------------------------------
\section{The ABM SDE}
%--------------------------------------------------------------

\begin{definition}[Arithmetic Brownian motion]\label{def:abm}
$X_t$ follows arithmetic Brownian motion with drift $\mu\in\R$ and
volatility $\sigma>0$ if it solves
\begin{equation}\label{eq:abmsde}
   \dd X_t = \mu\,\dt + \sigma\,\dW, \qquad X_0\in\R.
\end{equation}
In the notation of \texttt{ito\_foundations.tex} this is
$\drift(t,x)=\mu$, $\diff(t,x)=\sigma$ --- both constant, with no
$x$-dependence at all.
\end{definition}

\begin{remark}[Correspondence with the code]
\texttt{ABM::drift} returns the constant $\mu$ and
\texttt{ABM::diffusion} returns the constant $\sigma$, independent of
$x$ in both cases. Since neither coefficient depends on $x$, the
Lipschitz and linear-growth conditions of the existence--uniqueness
assumption in \texttt{ito\_foundations.tex} hold trivially --- this is
the simplest possible case of that assumption.
\end{remark}

%--------------------------------------------------------------
\section{Exact solution}
%--------------------------------------------------------------

Both GBM and OU needed a substitution to remove an $X_t$-dependence
from the coefficients before the SDE could be integrated: $\ln x$ for
GBM, the integrating factor $e^{\kappa t}$ for OU. Here there is
nothing to remove --- $\mu$ and $\sigma$ are already constants --- so
the general relation from \texttt{ito\_foundations.tex},
\[
   X_t = X_0 + \int_0^t \drift(s,X_s)\,\dd s + \int_0^t \diff(s,X_s)\,\dd W_s,
\]
integrates immediately once $\drift=\mu$ and $\diff=\sigma$ are
substituted.

\begin{theorem}[Exact terminal law of ABM]\label{thm:abm-exact}
Let $X_t$ solve \eqref{eq:abmsde}. Then for any $T>0$,
\begin{equation}\label{eq:abmexact}
   X_T = X_0 + \mu T + \sigma W_T,
\end{equation}
where $W_T\sim\mathcal N(0,T)$, so
\begin{equation}\label{eq:abmlaw}
   X_T \sim \mathcal N\big(X_0+\mu T,\ \sigma^2 T\big).
\end{equation}
\end{theorem}

\begin{proof}
Substituting $\drift(s,X_s)=\mu$ and $\diff(s,X_s)=\sigma$,
\[
   X_T = X_0 + \int_0^T \mu\,\dd s + \int_0^T \sigma\,\dd W_s
       = X_0 + \mu T + \sigma\int_0^T \dd W_s .
\]
The remaining integral is $\int_0^T \dd W_s = W_T - W_0 = W_T$,
giving \eqref{eq:abmexact}. Since $X_0+\mu T$ is a fixed number and
$W_T\sim\mathcal N(0,T)$, the affine transformation
$X_0+\mu T+\sigma W_T$ is itself Gaussian with mean $X_0+\mu T$ and
variance $\sigma^2\Var(W_T)=\sigma^2 T$, giving \eqref{eq:abmlaw}.
\end{proof}

\begin{remark}
Equation~\eqref{eq:abmlaw} is sampled with a single normal draw:
generate $Z\sim\mathcal N(0,1)$ and set
$X_T = X_0+\mu T+\sigma\sqrt T\,Z$. As with \texttt{GBM::sample\_terminal}
and \texttt{OU::sample\_terminal}, no time-stepping is needed --- this
is what \texttt{ABM::sample\_terminal} implements.
\end{remark}

%--------------------------------------------------------------
\section{Moments and the relation to OU}
%--------------------------------------------------------------

\Cref{thm:abm-exact} already gives the moments directly, with no
separate lemma needed:
\[
   \E[X_T] = X_0+\mu T, \qquad \Var(X_T) = \sigma^2 T.
\]
Unlike OU, the variance grows without bound as $T\to\infty$: there is
no mean-reverting pull holding the process in, so uncertainty simply
accumulates. Correspondingly, $\E[X_T]$ grows in a straight line
rather than curving toward a long-run level --- there is no $\theta$
here for it to curve toward.

\begin{remark}[ABM as a degenerate OU process]
Setting $\kappa=0$ in the OU drift $\kappa(\theta-X_t)$ removes it
entirely, leaving $\dd X_t=\sigma\,\dW$ --- driftless Brownian motion.
ABM generalises this by adding back a constant drift $\mu$ that is
\emph{not} tied to any mean-reversion mechanism. Equivalently: ABM is
what OU degenerates into once mean-reversion is switched off, plus an
unconstrained drift term OU's $\kappa=0$ case does not have on its
own.
\end{remark}

%--------------------------------------------------------------
\bibliographystyle{agsm}
\bibliography{refs}

\end{document}